\section{Conditional suppression: the logical benchmark and reconstruction}
\label{sec:annex-p19-conditional-suppression}

\begin{annexitem}{annex:p19.conditional-benchmark}{thm:p19-conditional-benchmark}{Truth preservation is independent of endorsement frequency}
\begin{proof}
Use the valuation clauses of Section~\ref{sec:prop-syntax-semantics-02}. If \(p\to q\)
and \(p\) are both true, \(q\) cannot be false, since that is exactly
the valuation making the conditional false. This proves MP. If
\(p\to q\) and \(\neg q\) are true, \(p\) cannot be true for the
same reason, proving MT. The valuation \(p=0,q=1\) satisfies both
the AC premises and the DA premises while falsifying each proposed
conclusion, proving their invalidity in this language.

If every valuation satisfying \(\Gamma\) satisfies \(q\), every
valuation satisfying \(\Gamma\cup\Delta\) also satisfies \(\Gamma\)
and hence \(q\). This is monotonicity, including when the enlarged
premise set has no satisfying valuation. Applying it to the MP premises
proves every claimed extension. In each resulting model \(q=1\), so
\(r\lor q=1\) as well. For the contravening and conjunctive forms,
\(p=1,q=1,r=0\) is a satisfying valuation, so their validity here
does not rest only on inconsistent premises. Their extra conditional
requires \(r=0\) once \(p=q=1\), in both cases.
None of these statements includes a frequency or a psychological
assumption; empirical endorsements are not their proof premises.
\end{proof}
\end{annexitem}

\begin{annexitem}{annex:p19.suppression-reconstruction}{thm:p19-suppression-reconstruction}{Countervaluations expose the changed interpretation}
\begin{proof}
The literal library premises include \(e\to s,e\), so MP gives
\(s\), regardless of the other conditional. For the revised set,
\((e,o,s)=(1,0,0)\) and \((1,0,1)\) both satisfy
\((e\land o)\to s\) and \(e\). Thus the set entails neither
\(s\) nor its negation. If \(o\) is added, \(e\land o\) is true
and MP gives \(s\) from the revised conditional.

The conjunction \((p\to q)\land(r\to q)\) is false exactly when
\(q=0\) and at least one of \(p,r\) is one. These are exactly the
falsifying valuations of \((p\lor r)\to q\), proving equivalence.
But \((p,r,q)=(1,0,0)\) makes \((p\land r)\to q\) true and
\(p\to q\) false. The conjunctive-antecedent interpretation is
therefore not an equivalence transformation of the original pair.

In the rain reconstruction the initial state has rows \((0,0,0)\)
and \((1,1,0)\), each satisfying \(p\to q\). Imposing the
categorical \(p\) retains only the second row, which satisfies \(q\).
The expanded final state has the complete table
\[
\begin{array}{ccc|ccccc}
 p&q&r&(p\land\neg r)\to q&r\to\neg q&p\to q&q&r\lor q\\\hline
 1&1&0&1&1&1&1&1\\
 1&0&1&1&1&0&0&1
\end{array}
\]
The table proves all the final-state claims. In particular the added
row is not a model of the original material conditional. This is a
change of the state over which the interpretation is evaluated, not
addition of a classical premise while retaining the original conditional
as true. Finally, if every valuation in a fixed nonempty state satisfies
both a material conditional and its antecedent, the already proved MP
argument applies to each valuation and makes the consequent true in all
of them. The reconstruction supplies no theorem about how many people
actually use either state or representation.
\end{proof}
\end{annexitem}
